\section{Introduction}
\label{sec:introduction}

Action chunking is a common interface for learned robot control: rather than
predicting only the next action, a policy predicts a temporally structured
sequence~\cite{zhao2023act,chi2025diffusionpolicy}.  The prediction does not by
itself specify how it should be deployed.  A controller may execute the whole
sequence open loop, or execute only a prefix before acquiring a new observation
and replanning.  The length of this prefix---the \emph{execution
horizon}---changes both the feedback rate and the frequency of transitions
between independently predicted chunks.

This choice is consequential.  Receding-horizon execution in Diffusion Policy
separates the action-prediction and action-execution horizons
\cite{chi2025diffusionpolicy}.  More recent work adapts chunk execution to
inference latency, prediction uncertainty, model attention, kinematic phase, or
learned value~\cite{black2025rtc,wang2026autohorizon,liang2026aac,
nie2026pace,feng2026dvac,zhao2026dehp,lei2026sparkvla}.  In the primary methods
we audited, the final decision is predominantly one prefix length or boundary
for the action vector as a whole.  There are important qualifications: AAC
separately measures translation, rotation, and gripper entropy, and PACE forms
arm-specific speed profiles.  Both nevertheless aggregate these signals before
choosing one synchronized boundary.

Robot action vectors themselves are structured.  A manipulation command can
combine end-effector motion and gripper actuation, or commands for multiple arms
and fingers.  Sharing one execution horizon forces every component to accept a
new prediction simultaneously.  We ask whether this synchronization is
necessary even when the policy was trained to predict the joint action sequence.

We isolate this question with a fixed, group-specific execution mechanism.  A
pretrained chunk policy remains frozen.  When one action group reaches its
horizon, the current observation produces a fresh full chunk; expired groups
accept their corresponding sequences, whereas non-expired groups retain their
existing sequences.  Consequently, an executed action can contain components
from different policy-query generations.  This mixed-generation composition is
both the enabling mechanism and a possible source of temporal inconsistency.

Static experiments with a frozen ACT checkpoint on LIBERO Object
\cite{liu2023libero} reveal systematic group-dependent sensitivity.  On one task
we evaluate all 25 arm/gripper assignments from horizons
$\{1,2,4,8,16\}$ over all 50 official initial states.  The best evaluated pair
improves a global control at matched policy-query rate, and all ten symmetric
horizon swaps favor the assignment with the longer gripper horizon.  Across the
other nine object tasks, off-diagonal pairs are repeatedly competitive and six
tasks exhibit a strict improvement over a global point without a higher query
rate.  These observations do not show that the gripper should generally update
more slowly, nor do they establish a dynamic scheduler.  They show that assigning
the same two horizon values to different physical groups can change execution
outcomes while preserving policy-query cadence.

This draft makes four evidence-bounded contributions:
\begin{itemize}
    \item a controlled execution formulation that permits physically meaningful
    action groups to use different fixed horizons without changing the
    pretrained chunk policy;
    \item a complete $5\!\times\!5$ arm/gripper horizon landscape on all 50
    official states of one LIBERO Object task, with diagonal controls and query
    accounting;
    \item a coarse cross-task diagnostic over the remaining nine LIBERO Object
    tasks, separating deployable static settings from retrospective per-task
    oracle-style comparisons; and
    \item paired, budget-controlled evidence that swapping or separating group
    horizons can change success under essentially identical policy-query rates.
\end{itemize}
